\documentclass[10pt,a4paper]{amsart}
\usepackage[margin=19mm]{geometry}
\usepackage{amsmath,amssymb,mathtools}
\usepackage[scr=rsfs,cal=euler]{mathalfa}
\usepackage{xfrac}
\usepackage[unicode,hidelinks,bookmarks=false]{hyperref}
\newcommand{\ie}{i.e.\ }
\newcommand{\cf}{cf.\ }
\newcommand{\resp}{respectively\ }
\newcommand{\Subsection}{\subsection}
\providecommand{\nobreakdash}{\nobreak\hbox{-}}
\setlength{\emergencystretch}{2em}
\multlinegap0pt
\usepackage{tikz}
\usepackage{enumerate}
\usepackage{amssymb}
\newtheorem{theorem}{Theorem}
\newcommand{\alg}[1]{{\textbf{\upshape #1}}}  %
\newcommand{\back}{\backslash}
\newcommand{\betafil}{\mathbf{\beta Fil}}
\newcommand{\dbetafil}{{\rm \beta Fil}}
\newcommand{\dfil}{{\rm Fil}}
\newcommand{\fil}{\mathbf{Fil}}
\newcommand{\gell}{\gamma^\ell}
\newcommand{\gr}{\gamma^r}
\newcommand{\sell}{\sigma^\ell}
\newcommand{\sr}{\sigma^r}
\newcommand{\sse}{\subseteq}
\title{Blockwise Gluings And Amalgamation Failures\\ in Integral Residuated Lattices}
\author{Valeria Giustarini, Sara Ugolini}
\date{Source: arXiv:2408.17400v3 (CC BY 4.0)}
\begin{document}
\maketitle

\noindent\emph{Source and licence.} Blockwise Gluings And Amalgamation Failures\\ in Integral Residuated Lattices. Version arXiv:2408.17400v3 (CC BY 4.0), distributed by arXiv under the Creative Commons Attribution 4.0 International licence, http://creativecommons.org/licenses/by/4.0/, as recorded in the arXiv licence metadata for this exact version and shown on the abstract page https://arxiv.org/abs/2408.17400v3. Full licence notice: \texttt{licenses/arxiv-2408.17400-CC-BY-4.0.txt}.\par\smallskip

\noindent\textit{Editorial scope.} Counted unit: the article's theorem filters (source0.tex lines 778--784). The article's complete original proof of it is delivered with it (source0.tex lines 785--829, one \texttt{proof} environment of 45 lines). No mathematics is written, repaired or re-derived: the statement and the proof are the article's own lines, in the article's own notation, and no retained line is substituted or re-flowed.  No further source-local context is needed: every cross-reference of every retained line resolves inside the two delivered spans.  Nothing was omitted from any retained span, so the package carries no omission record.  No retained line contains a citation, so the wrapper carries no bibliography and no bibliography entry.  No retained line contains a figure environment, an image-inclusion command or drawing code, so no image is delivered and the package declares no asset dependency.  Editorial formatting only, in addition: an AMS \texttt{amsart} preamble that restates the article's own declarations and macros used by the retained lines, namely the environments \texttt{theorem} on separate counters, together with the standard packages those lines need.  The retained environments print in this document as Theorem 1 in order of appearance, whereas the article prints the same items with the numbering of its own full document.  The complete native source of the article as distributed in the arXiv e-print for this exact version is bundled unchanged as \texttt{sources/164526/source0.tex}.
\begin{theorem}\label{theorem: description of filters}
	Let $\alg B \boxplus_\beta \alg C$ with $\beta=(\sell, \sr, \gell, \sr)$. Then 
	\begin{enumerate}
		\item if $\sell=\sr$, $\fil(\alg B \boxplus_\beta \alg C)\cong \fil(\alg C) \oplus (\betafil(\alg B) -\{1\})$, where $\oplus$ is the ordinal sum of posets;
		\item if $\sell \neq \sr$, $\fil(\alg B \boxplus_\beta \alg C)\cong\betafil(\alg B)$.
	\end{enumerate}
\end{theorem}

\begin{proof}
	(1)  Let us consider the map $\varphi: \fil(\alg B \boxplus_\beta \alg C) \to \fil(\alg C) \oplus (\betafil(\alg B)-\{1\})$ defined by:
\begin{align*}
	\varphi(F) = &\left\{\begin{array}{l l}
F &\mbox{ if } F \cap B = \{1\}, \vspace{0.1cm}\\
F \cap B &\mbox{ otherwise.} \vspace{0.1cm}
\end{array}\right.
\end{align*}
The map $\varphi$ is well defined. Indeed, let $F \in \dfil(\alg B \boxplus_\beta \alg C)$, then if $F \cap B = \{1\}$, clearly $F \in \dfil(\alg C)$. 
	  If instead there exists an $ x \in F \cap  B$ with $x \neq 1$, then notice that $C \subseteq F$, since $F$ needs to be closed upwards, and it is easy to see that $F \cap  B \in \dbetafil(\alg B)-\{1\}$. 
	  
	  We prove that $\varphi$ is the desired lattice isomorphism. It is a straightforward consequence of the definition that $\phi$ is injective and preserves the order. To prove surjectivity, let $F \in \dfil(\alg C) \oplus (\dbetafil(\alg B)-\{1\})$, then either $F \in \dfil(\alg C)$ or $F \in \dbetafil(\alg B)-\{1\}$. 
	  
	  If $F \in \dfil(\alg C)$, we show that $F$ is also a congruence filter of $\alg B \boxplus_\beta \alg C$, and hence $\varphi(F) = F$. Since $\alg C$ is a subalgebra of the gluing, it suffices to check that $F$ is closed under conjugates with respect to elements of $\alg B$. Suppose there is $x \in F, x \neq 1$, and let $b \in B-\{1\}$. Then $b \backslash xb = b \backslash \sell(b) = c_{\alg C}$ by the definition of the operations, since $\sr(b) = \sell(b)$ by hypothesis. Similarly, $bx/b = c_{\alg C}$. Notice that $x \leq c_{\alg C}$ hence $c_{\alg C} \in F$. Thus, $F$ is a filter of the gluing as desired. 
	  
	  Suppose now $F \in \dbetafil(\alg B)$, $F \neq\{1\}$, and let $F' = F \cup C$. We prove that $F'$ is a congruence filter of the gluing. It suffices to check closure under conjugates involving elements of $\alg B$, since $F'$ is clearly closed under products and upwards, and $\alg C$ is a subalgebra of the gluing. We need to consider the following cases:
	  \begin{itemize}
	  	\item If $x \in C, b \in B$, then $bx/b = \sr(b) / b \in C\sse F'$ by the definition of the operations, since $\sell(b) = \sr(b)$. Similarly $b \back xb \in C \sse F'$.
	  	\item  Let now $x \in F, b \in B$, then $bx/b$ and $b \backslash xb$ are either in $B$, hence in $F$ given that $F$ is a partial filter, or in $C$, hence in $F'$.
	  	\item Consider then $x \in F, c \in C$. Then $cx / c = \sell(x) / c = \gr\sell(x)$. Now, since $\sr(x) = \sell(x)$ by hypothesis, we get $x \leq  \gr\sell(x) \in F\sse F'$, using that $(\sr, \gr)$ is a residuated pair. 
	  \end{itemize}
	   Hence $F'$ is a filter of $\alg B \boxplus_\beta \alg C$, and clearly $\varphi(F') = F$. We have shown that $\varphi$ is an order-preserving bijection. To see that it is a lattice isomorphism, it suffices to observe that, given what we have shown above, the map $\varphi': \fil(\alg C) \oplus (\betafil(\alg B)-\{1\}) \to  \fil(\alg B \boxplus_\beta \alg C)$ defined by:
	    \begin{align*}
	\varphi'(F) = &\left\{\begin{array}{l l}
F &\mbox{ if } F \in \dfil(\alg C), \vspace{0.1cm}\\
F \cup C &\mbox{ if } F \in \dbetafil(\alg B) \vspace{0.1cm}
\end{array}\right.
\end{align*}
	   is the inverse of $\varphi$, and it is also order preserving.
	   This proves that $\varphi$ is the desired lattice isomorphism.
	 
	 (2) Suppose now that $\sell \neq \sr$. First, we notice that every nontrivial filter $F \in  \dfil(\alg B \boxplus_\beta \alg C)$ intersects $B - \{1\}$. Indeed, let $ x\in F-\{1\}$. Since  $F$ is closed under conjugation, for any $y \in B$ and $x \in F $, $y x / y$ and $y \backslash x y $ are in $F$. Recalling that $\sell \neq \sr$, there exists $b \in B$ such that $\sell(b)\neq \sr(b)$, i.e. either $\sell(b)\not \leq \sr(b)$ or $\sr(b)\not \leq \sell(b)$. If $\sell(b)\not \leq \sr(b)$, then, by definition of the operations, $b x / b=\sr(b)/b \in B -\{1\}$. Similarly, if $\sr(b)\not \leq \sell(b)$, then $b \backslash x b=b \backslash \sell(b) \in B-\{1\}$. 

Let us then consider the map $\psi: \fil(\alg B \boxplus_\beta \alg C) \to \betafil(\alg B)$ sending $F \mapsto F \cap B$. To see that $\psi$ is indeed a map from the filters of the gluing to the $\beta$-filters of the lower block $(\alg B, \beta)$, consider any nontrivial $F \in \dfil(\alg B \boxplus_\beta \alg C)$. It is clear that $F \cap B$ is a partial filter of $\alg B$. To see that it is a $\beta$-filter, notice first that, as argued above, any $c \in C$ is in $F$. Thus, by closure under conjugates, for every $b \in B$, both $b \backslash cb = b \backslash \sell(b)$ and $bc / b = \sr(b) / b$ are in $F$. Therefore, if $\sell(b) \not\leq\sr(b)$ we have that $bc / b = \sr(b) /b \in F \cap B$, and similarly, for any $b \in B$ such that $\sr(b) \not\leq \sell(b)$ we have that $b \backslash \sell(b) \in F \cap B$. Hence, $F \cap B$ is a $\beta$-filter.

 We move to show that $\psi$ is the desired lattice isomorphism. As above, $\psi$ is clearly injective and order preserving. To prove surjectivity, Note first that $\psi(\{1\}) = \{1\}$. Now
 let $F \in \dbetafil(\alg B) - \{1\}$, and consider $F' = F \cup C$. We prove that $F' \in \dfil(\alg B \boxplus_\beta \alg C)$, by showing that it is closed under conjugates. We need to consider the following cases:
 \begin{itemize}
  	\item Let first $x \in C, b \in B$, then $bx /b = \sr(b) / b$. If $\sell(y) \leq \sr(y)$ then $\sr(b) / b \in C \sse F'$. If instead $\sell(y) \not\leq \sr(y)$, then $\sr(b) / b \in F$ since $F$ is a $\beta$-filter. The proof that $b \backslash xb \in F$ is similar. 
 	\item Let now $x \in F, b \in B$, then $bx /b$ and $b \backslash xb$ are either in $B$, hence in $F \sse F'$ because $F$ is closed under existing conjugates, or in $C$, hence in $F'$.
 	\item Finally, let $x \in F, c \in C$. Then $cx/c = \sell(x)/c \geq \sell(x) \geq x^2 \in F$, and then $cx/c \in F$. Similarly $c \backslash xc \in F$. 
 \end{itemize}

Thus $F'$ is a congruence filter of the gluing, and since $\psi(F') = F$, we proved surjectivity of $\psi$. It also follows from what we have shown that the map $\psi': \betafil(\alg B) \to \fil(\alg B \boxplus_\beta \alg C)$ defined by $\psi'(\{1\}) = \{1\}$ and $\psi'(F) = F \cup C$ is the inverse of $\psi$, and $\psi'$ is also clearly order-preserving. Therefore, $\psi$ provides the lattice isomorphism showing $\fil(\alg B \boxplus_\beta \alg C)\cong\betafil(\alg B)$.
\end{proof}

\end{document}
